\chapter{Penutup}
\label{ch5}

\section{Kesimpulan}

Berdasarkan hasil implementasi dan pengujian yang telah dibahas pada Bab~\ref{ch4}, kesimpulan penelitian ini disusun sesuai dengan tiga rumusan masalah yang telah ditetapkan pada Bab~\ref{ch1}:

\begin{enumerate}
	\item Mengatasi Pencatatan Manual tanpa Jejak Audit. Empat \textit{smart contract} Lakomi (Token, Vault, Govern, Loans) telah berhasil dirancang dan diterapkan pada jaringan DChain (Chain ID 17845) dengan arsitektur \textit{dual-track} yang memisahkan hak suara satu-anggota-satu-suara dari insentif kontribusi berbasis LTV, serta menerapkan \textit{Role-Based Access Control} delapan peran yang menjamin pemisahan kekuasaan. Sistem mencatat seluruh transaksi koperasi, pendaftaran anggota (93.344 gas), pembayaran simpanan (127.526 gas), pemungutan suara (110.228 gas), dan distribusi SHU (\textasciitilde200.000 gas), secara permanen pada \textit{distributed ledger} DChain yang dapat diaudit secara independen oleh setiap anggota.

	\item Mengatasi Ketergantungan pada Kepatuhan Sukarela. Sebanyak 26 ketentuan dari UU 25/1992, PP 7/2021, dan Permenkop 8/2023 telah berhasil dipetakan ke dalam fungsi \textit{smart contract} yang spesifik dan terverifikasi. Seluruh 26 skenario pengujian fungsional berbasis Hardhat menunjukkan hasil Lulus (100\%), mencakup pengujian \textit{happy path}, \textit{error path}, dan verifikasi kontrol akses. Kepatuhan kini ditegakkan oleh kode yang \textit{self-executing}, bukan oleh kepatuhan sukarela pengurus.

	\item Menyediakan Data Perbandingan Sistem Koperasi. Analisis gas pada DChain menunjukkan total konsumsi \textit{deployment} sebesar 12,2 juta gas untuk empat kontrak dengan biaya operasional per transaksi di bawah 0,001 DKT. Satu siklus pinjaman penuh mengonsumsi sekitar 620.000 gas. Perbandingan 9 dimensi terhadap koperasi konvensional dan DAO \textit{token-weighted} menegaskan bahwa Lakomi unggul dalam tiga aspek utama: transparansi \textit{ledger} publik, demokrasi satu-anggota-satu-suara yang dijamin \textit{smart contract}, dan otomatisasi distribusi SHU enam kategori secara on-chain.
\end{enumerate}

\section{Saran}

Berdasarkan batasan penelitian dan temuan selama pembahasan, saran untuk pengembangan selanjutnya dikelompokkan sebagai berikut:

\subsection{Berdasarkan Batasan Penelitian}

\begin{enumerate}
	\item Menerapkan sistem pada DChain \textit{mainnet} dengan anggota riil untuk memvalidasi kinerja, skalabilitas, dan konsumsi gas dalam kondisi produksi.

	\item Mengintegrasikan identitas digital (KTP elektronik atau \textit{Self-Sovereign Identity}) untuk memperkuat verifikasi keanggotaan dan mendukung kepatuhan KYC.

	\item Mengembangkan antarmuka yang ramah pengguna dengan panduan interaktif dan dukungan multi-bahasa untuk menjangkau anggota dengan literasi digital terbatas.

	\item Melakukan audit keamanan formal oleh pihak ketiga serta \textit{formal verification} terhadap properti kritis seperti \textit{invariant} saldo dan integritas kontrol akses.

	\item Memperluas cakupan regulasi ke Peraturan OJK, peraturan perpajakan SHU, dan integrasi pelaporan Kementerian Koperasi dan UKM.
\end{enumerate}

\subsection{Berdasarkan Temuan Pembahasan}

\begin{enumerate}
	\item Mengevaluasi efektivitas parameter tata kelola (kuorum 67\%, periode voting 7 hari, \textit{timelock} 24 jam) melalui studi longitudinal pada koperasi yang beroperasi penuh secara on-chain.

	\item Mengeksplorasi interoperabilitas \textit{cross-cooperative lending} yang memungkinkan anggota antar-koperasi saling meminjam melalui \textit{smart contract} \textit{cross-chain}.

	\item Mengintegrasikan \textit{zero-knowledge proofs} untuk meningkatkan privasi data anggota sambil mempertahankan auditabilitas publik atas kepatuhan kolektif.

	\item Menguji kelayakan \textit{decentralized identity} (DID) dan \textit{verifiable credentials} untuk verifikasi kelayakan kredit anggota tanpa mengungkapkan riwayat keuangan lengkap.
\end{enumerate}
